\section*{RESUMEN}
\addcontentsline{toc}{section}{RESUMEN}

Se desarrolló un sistema de gestión de acceso vehicular para un parqueadero de
Oruro con el fin de sustituir el registro manual de ingresos y salidas por un
proceso automatizado de visión artificial y aprendizaje profundo. El sistema se
apoyó en las cámaras existentes del recinto e integró un conjunto de datos propio,
modelos de detección y clasificación, un reconocedor de matrículas y una plataforma
web con base de datos, cola de tareas y notificación en tiempo real. Se verificaron
640 matrículas reales y se generaron 360,000 muestras sintéticas. La detección de
vehículos alcanzó un mAP@50 de 0.901; la de placas, una exhaustividad de 0.9714; la
clasificación de tipo, un F1 de 0.8887; y la de color, un F1 macro de 0.8124. El
reconocimiento de matrículas obtuvo 0.851 de exactitud por carácter y 52.6\,\% de
coincidencia exacta por recorte, ambos limitados por la resolución de captura. Dos
criterios quedaron por debajo del umbral fijado (el F1 de color y la exactitud
por carácter), la determinación de eventos quedó inconclusa por tamaño de
muestra y la latencia por evento no llegó a medirse; por ello se rechaza la
hipótesis de investigación y se acepta la hipótesis nula conforme al criterio
definido de antemano. La estimación económica arrojó 62 puntos de función y un
costo total de 33,170 bolivianos. El sistema está construido e integrado por
componentes, pero no satisface el conjunto de criterios de desempeño en su estado
actual.

\vspace{1em}
\noindent\textbf{Palabras clave:} visión artificial, aprendizaje profundo,
reconocimiento de matrículas, gestión de acceso, parqueadero.
